\begin{table}[!htbp]
\centering
\caption{Data sources, survey timing, and sampling frameworks.}
\label{tab:s1_cohort_sources}
\fontsize{9}{11}\selectfont
\setlength{\tabcolsep}{3.5pt}
\renewcommand{\arraystretch}{1.18}

\parbox{\textwidth}{\raggedright\textbf{A. Cohort data sources}}\par
\vspace{0.2em}
\begin{tabularx}{\textwidth}{@{}L{1.15cm}L{1.75cm}L{5.05cm}L{1.75cm}Y@{}}
\toprule
Cohort & Country & Harmonized data source & Available waves & Survey period \\
\midrule
HRS & United States & RAND HRS Longitudinal File 2022, Version 1 & W1--W16 & 1992--2022 waves; some interviews continued into 2023 \\
ELSA & England & Harmonized ELSA, Version H & W1--W10 & 2002--2023 \\
KLoSA & South Korea & Gateway Harmonized KLoSA, Version F & W1--W9 & 2006--2022 \\
\bottomrule
\end{tabularx}

\vspace{0.85em}

\parbox{\textwidth}{\raggedright\textbf{B. Survey-wave timing windows}}\par
\vspace{0.2em}
\begin{tabularx}{\textwidth}{@{}L{1.15cm}L{3.10cm}L{3.10cm}Y@{}}
\toprule
Cohort & Earliest supported four-wave window & Latest supported four-wave window & Timing note \\
\midrule
HRS & W2--W5 & W13--W16 & The HRS FI starts at W3; the earliest index-shock wave is W3. \\
ELSA & W1--W4 & W7--W10 & W10 fieldwork occurred in 2021--2023 after COVID-19 disruption. \\
KLoSA & W1--W4 & W6--W9 & Waves were generally conducted at approximately 2-year intervals. \\
\bottomrule
\end{tabularx}

\vspace{0.85em}
\parbox{\textwidth}{\raggedright\textbf{C. Original survey sampling frameworks and recruitment}}\par
\vspace{0.2em}
\begin{tabularx}{\textwidth}{@{}L{1.15cm}L{4.35cm}L{5.15cm}Y@{}}
\toprule
Cohort & Original target population and sampling frame & Sampling and recruitment features & Key coverage limitations \\
\midrule
HRS & Original 1992 cohort: adults born in 1931--1941, living in households in the contiguous United States. & Stratified multistage area probability sampling. Spouses or partners included regardless of age. Other birth cohorts were integrated, with later replenishment. & Community-based recruitment initially excluded institutional residents. Participants entering institutions could be followed subsequently. \\
ELSA & Adults aged 50 years and over in private households in England, recruited through previous Health Survey for England (HSE) samples. & HSE-based stratified multistage sampling. Eligible household members and partners, including younger partners, were recruited. Refreshment samples were added at later waves. & England, not the whole United Kingdom. The original household frame excluded institutional residents. Recruitment depended on earlier HSE participation and recontact eligibility. \\
KLoSA & Adults aged 45 years and over in private households in South Korea, using 2005 Census enumeration districts. & Stratified multistage household probability sampling. All age-eligible household members were eligible, but younger spouses were not interviewed. A refreshment sample was added in 2014. & Jeju and island or facility enumeration districts were excluded from the original frame. New spouses outside the baseline household sample were not recruited. \\
\bottomrule
\end{tabularx}

\vspace{0.25em}
\parbox{\textwidth}{\raggedright\fontsize{9}{11}\selectfont\textit{Notes:} Each four-wave sequence comprises the pre-shock ($t-1$), index-shock ($t$), landmark ($t+1$), and next scheduled follow-up ($t+2$) waves. The earliest and latest sequences shown are those available within each cohort's survey coverage; they do not describe the distribution of participants in the final analytic sample or imply that every participant was interviewed at all four waves. Published survey response and retention indicators differ in their reference populations, denominators, and treatment of deaths and are therefore not directly comparable across cohorts. Study-specific selection and observed frailty follow-up are reported separately in Supplementary Figure S1 and Tables S5--S6. Harmonised definitions support conceptual alignment, not complete measurement equivalence (Supplementary Tables S2--S4). HRS, Health and Retirement Study; ELSA, English Longitudinal Study of Ageing; KLoSA, Korean Longitudinal Study of Aging.}
\vspace{0.25em}
\parbox{\textwidth}{\raggedright\fontsize{9}{11}\selectfont\textit{Sources for Panel C and participation definitions:} \href{https://doi.org/10.1093/ije/dyu067}{Sonnega et al. (2014), HRS cohort profile}; \href{https://www.niaaa.nih.gov/research/alcohol-epidemiologic-data-directory-2022}{NIAAA survey directory (2022), HRS entry}; \href{https://hrs.isr.umich.edu/sites/default/files/biblio/ResponseRates_2017.pdf}{HRS response-rate documentation (2017)}; \href{https://doi.org/10.25553/h52b-8869}{Harmonized ELSA Version H (2025)}, Sections 1--1.2, and ELSA Wave 1 and Wave 10 technical reports (2007, 2024); \href{https://doi.org/10.25553/0gba-4r74}{Harmonized KLoSA Version F (2025)}, Sections 1--1.2; \href{https://survey.keis.or.kr/klosa/klosa03.jsp}{KEIS sampling documentation} and \href{https://www.kostat.go.kr/boardDownload.es?bid=12030&list_no=434382&seq=1}{KLoSA statistical information report (updated 2024), Table IV-5}.}
\end{table}
